\documentclass[12pt,a4paper]{article}
\usepackage[utf8]{inputenc}
\usepackage[T1]{fontenc}
\usepackage{geometry}
\geometry{margin=2.5cm}
\usepackage{hyperref}
\usepackage{booktabs}
\usepackage{enumitem}
\usepackage{xcolor}
\usepackage{tcolorbox}

\definecolor{addedcolor}{RGB}{0,100,0}

\newtcolorbox{addedsource}{
  colback=green!5!white,
  colframe=green!50!black,
  title={\textbf{Additional Source (not in original proposal)}},
  fonttitle=\small
}

\title{\textbf{ShinySwarm: Literature Review Summary}\\[0.5em]
\large For Midterm Presentation Preparation}
\author{Iva Gunjaca\\
Master of Software Engineering, University of Amsterdam}
\date{\today}

\begin{document}
\maketitle
\tableofcontents
\newpage

% ============================================================
\section{The Reproducibility Crisis and R's Role in Science}
% ============================================================

Nature's survey of 1,576 researchers found that more than 70\% of experiments could not be reproduced by independent peers, and more than 50\% could not be reproduced by the original authors themselves \cite{baker2016}. Peng \cite{peng2011} argues that lack of reproducibility stems from poorly done data analysis that is not publicly available --- once code and data are open, errors become identifiable.

R \cite{rcore2025} has become a mainstay for statistical analysis and data visualisation in the scientific community. However, it requires programming literacy that remains a barrier for domain experts. The Shiny package \cite{chang2022} addresses this by enabling interactive web applications using only R, without HTML, CSS, or JavaScript knowledge, broadening accessibility for non-programmers \cite{hartgerink2017}.

\begin{addedsource}
Stodden, V., McNutt, M., Bailey, D.H. et al. ``Enhancing reproducibility for computational methods.'' \textit{Science}, 354(6317), 2016. --- Argues that containerisation of computational environments is a key enabler for reproducibility. Directly supports the Docker-based approach used in ShinySwarm.
\end{addedsource}

% ============================================================
\section{Collaboration Models for Shiny Applications}
% ============================================================

The collaboration framework is grounded in Churchill et al.'s \cite{churchill2001} definition of Collaborative Virtual Environments (CVEs): ``distributed virtual reality systems providing potentially limitless virtual spaces where users can engage in information sharing.''

Interactions with Shiny web applications can be classified as \textbf{read} (viewing outputs) and \textbf{write} (changing inputs/state), leading to five collaboration modes:

\begin{itemize}[nosep]
  \item \textbf{Read--Read} --- multiple users observe simultaneously (natively supported by Shiny)
  \item \textbf{Write--Write} --- multiple users alter state simultaneously (explored but problematic)
  \item \textbf{Read--Write} --- hybrid where some users read, some write (the OWNER/EDITOR/VIEWER model in ShinySwarm)
  \item \textbf{Individual interaction} --- no interference between users
  \item \textbf{State sharing} --- transferring current application state between users
\end{itemize}

\begin{addedsource}
\begin{itemize}[nosep]
  \item Ellis, C.A. and Gibbs, S.J. ``Concurrency control in groupware systems.'' \textit{ACM SIGMOD Record}, 18(2), 1989. --- The foundational taxonomy of collaboration modes (same-time/same-place matrix). The read/write classification maps directly onto their framework.
  \item Sun, C. and Ellis, C. ``Operational transformation in real-time group editors: issues, algorithms, and achievements.'' \textit{Proceedings of ACM CSCW}, 1998. --- The theory behind Google Docs-style collaboration. Relevant because \texttt{shiny.collections} attempted this approach and failed.
\end{itemize}
\end{addedsource}

% ============================================================
\section{The State Sharing Problem}
% ============================================================

The \texttt{shinyURL} package \cite{aerts2015} and Shiny bookmarking \cite{shinybookmark} encode input values into URLs for state sharing. The key limitation: \textbf{this requires double computation} --- the receiving user must recompute all outputs from the shared inputs.

The solution adopted in ShinySwarm: transfer \textbf{both inputs and outputs}, so inputs provide understanding of methodology while outputs prevent redundant computation. This became the design principle behind the Redis/Kafka state relay architecture.

The \texttt{shiny.collections} package \cite{appsilon2017} attempted Google Docs-style live collaboration using RethinkDB \cite{rethinkdb} as a reactive database. It was presented at UseR! 2017 but is now \textbf{deprecated} because RethinkDB lost R support. The shared-database approach also introduces the anti-pattern identified by Richardson \cite{richardson_shared_db}: concurrent read/write operations risk data corruption, and a bug in one service can affect all others.

There has been no other publicly documented exploration of write--write or read--write collaboration for Shiny applications.

% ============================================================
\section{Shiny Deployment Platforms}
% ============================================================

Three deployment approaches exist for Shiny applications:

\begin{table}[h]
\centering
\small
\begin{tabular}{@{}lllll@{}}
\toprule
\textbf{Platform} & \textbf{Architecture} & \textbf{Scaling} & \textbf{Auth} & \textbf{Containers} \\
\midrule
shinyapps.io \cite{shinyappsio} & Multi-tenant cloud & Auto-scaling R processes & Paid tiers & Behind the scenes \\
Shiny Server OSS \cite{shinyserver} & Self-hosted Linux & Manual, no LB & None built-in & None native \\
ShinyProxy \cite{shinyproxy} & Container-per-session & Horizontal (K8s/Swarm) & LDAP/AD & Core architecture \\
\bottomrule
\end{tabular}
\caption{Comparison of existing Shiny deployment platforms}
\end{table}

ShinyProxy is the closest existing solution to ShinySwarm --- it uses Spring Boot on the server side and launches isolated Docker containers per user session. However, \textbf{none of these platforms support real-time multi-user collaboration on the same application instance}.

\begin{addedsource}
\begin{itemize}[nosep]
  \item Fay, C., Rochette, S., Guyader, V., Girard, C. \textit{Engineering Production-Grade Shiny Apps.} Chapman \& Hall/CRC, 2021. --- The \texttt{golem} framework book. Discusses Shiny app architecture at scale, relevant context for why monolithic Shiny apps hit limits.
  \item OpenAnalytics. ``ShinyProxy Architecture.'' shinyproxy.io/documentation. --- Details the Spring Boot + Docker container-per-session model that directly inspired the ShinySwarm architecture.
\end{itemize}
\end{addedsource}

% ============================================================
\section{Microservice Architecture and Shiny}
% ============================================================

R Shiny apps are typically monolithic. Only \textbf{two prior attempts} at microservice-based Shiny exist in the literature:

\begin{enumerate}[nosep]
  \item \textbf{``Shiny Microservices Demo''} --- useR! 2018 Conference \cite{brokamp2018}. GitHub-only, minimal documentation. Uses an API gateway for inter-service communication.
  \item \textbf{``Turning our monolithic Shiny app into a microservices-based structure''} --- Juan Cruz Rodriguez \cite{rodriguez2021}. Uses Plumber \cite{plumber2025} to expose R functions as REST APIs with roxygen2-like comments.
\end{enumerate}

Both use API gateway + Plumber for communication. Neither implements real-time collaboration.

\begin{addedsource}
\begin{itemize}[nosep]
  \item Lewis, J. and Fowler, M. ``Microservices: a definition of this new architectural term.'' martinfowler.com, 2014. --- The canonical definition: ``suites of small services, each running in its own process and communicating with lightweight mechanisms, often an HTTP resource API.'' Directly frames the REST vs.\ Kafka comparison.
  \item Newman, S. \textit{Building Microservices.} O'Reilly, 2nd edition, 2021. --- Industry standard reference. Chapter on inter-service communication covers synchronous (REST) vs.\ asynchronous (messaging) patterns --- the exact comparison being made.
  \item Richardson, C. \textit{Microservices Patterns.} Manning, 2018. --- Defines the API Gateway pattern, Database-per-Service pattern, and Saga pattern. The Spring Boot backend implements the API Gateway pattern; the Redis/Kafka state stores implement Database-per-Service.
\end{itemize}
\end{addedsource}

% ============================================================
\section{Apache Kafka and R Integration}
% ============================================================

Apache Kafka \cite{kafka_docs} is a distributed event streaming platform originally created at LinkedIn. It uses publish-subscribe messaging with topics, partitions, and consumer groups.

The R--Kafka integration history:

\begin{enumerate}[nosep]
  \item \textbf{\texttt{rkafka} package} \cite{gupta2022} --- Removed from CRAN, only available on GitHub, not updated since 2022. The author's own Proof-of-Concept using \texttt{rkafka} showed \textbf{21-second end-to-end latency} due to dataframe-to-JSON serialisation taking 5--8 seconds per message.
  \item \textbf{Neff's MongoDB approach} \cite{neff2019} --- Store Kafka messages in MongoDB, read from R. Only supports consumption, not production.
  \item \textbf{\texttt{kafka} package by INWT} \cite{neudecker2024} --- New R package funded by R Consortium, built on \texttt{librdkafka}. Provides \texttt{Producer\$new()} and \texttt{Consumer\$new()} classes. The final Kafka architecture uses this package, which resolved the latency issue.
\end{enumerate}

\begin{addedsource}
\begin{itemize}[nosep]
  \item Kreps, J. ``The Log: What every software engineer should know about real-time data's unifying abstraction.'' LinkedIn Engineering Blog, 2013. --- Jay Kreps (Kafka co-creator) explains the log abstraction that underpins Kafka. Directly relevant to understanding why Kafka's ordered log is suitable for state synchronisation in collaborative apps.
  \item Narkhede, N., Shapira, G., Palino, T. \textit{Kafka: The Definitive Guide.} O'Reilly, 2nd edition, 2021. --- Covers consumer groups, partitioning, and exactly-once semantics. The use of message keys for session routing is a direct application of Kafka's partitioning model.
  \item INWTlab. ``r-kafka: R Client for Apache Kafka.'' GitHub, github.com/INWTlab/r-kafka. --- Funded by R Consortium. Documents the \texttt{Producer\$new()} / \texttt{Consumer\$new()} API used throughout the Kafka architecture.
\end{itemize}
\end{addedsource}

% ============================================================
\section{Data Persistence in Shiny}
% ============================================================

Native Shiny provides no built-in data persistence. Attali \cite{attali2016} documents seven recommended methods: local files (CSV, RDS), SQLite, relational databases (MySQL, PostgreSQL), MongoDB \cite{mongodb}, Google Sheets, Dropbox, and Amazon S3.

ShinySwarm uses \textbf{PostgreSQL} for persistent checkpoint storage (the \texttt{SavedState} entity) and \textbf{Redis} (REST architecture) or \textbf{Kafka topics + ConcurrentHashMap} (Kafka architecture) for live session state.

\begin{addedsource}
Carlson, J.L. \textit{Redis in Action.} Manning, 2013. --- Covers Redis as a real-time cache and pub/sub system. The REST architecture uses Redis as the ``State Vault'' --- the central source of truth that all R Shiny frontends poll.
\end{addedsource}

% ============================================================
\section{Containerisation and Reproducibility}
% ============================================================

Docker \cite{docker_docs} encapsulates the entire computational environment --- code, dependencies, system libraries --- ensuring analyses can be reliably re-executed across systems. Containerisation serves as the bridge between reproducibility and collaboration in this work.

\begin{addedsource}
\begin{itemize}[nosep]
  \item Boettiger, C. ``An introduction to Docker for reproducible research.'' \textit{ACM SIGOPS Operating Systems Review}, 49(1), 2015. --- The key paper on Docker for scientific reproducibility. Argues that containerisation solves the ``works on my machine'' problem for computational research. Highly cited (1000+), directly supports the thesis motivation.
  \item Merkel, D. ``Docker: lightweight Linux containers for consistent development and deployment.'' \textit{Linux Journal}, 239, 2014. --- Early technical overview of Docker's architecture (cgroups, namespaces, layered filesystems).
  \item Burns, B. et al. ``Borg, Omega, and Kubernetes.'' \textit{ACM Queue}, 14(1), 2016. --- Google's paper on container orchestration evolution. Relevant to the Rancher/Kubernetes deployment target.
\end{itemize}
\end{addedsource}

% ============================================================
\section{Presentation Suggestions}
% ============================================================

For a 30-minute midterm presentation, the literature review should occupy \textbf{2--3 slides maximum}:

\begin{enumerate}
  \item \textbf{Slide: ``The Gap''} --- Reproducibility crisis (Baker 2016) + Shiny is single-user + existing collaboration attempts failed (\texttt{shiny.collections} deprecated) + no microservice Shiny exists with real-time collaboration. This is the motivation.
  
  \item \textbf{Slide: ``Existing Solutions \& Their Limits''} --- Table comparing shinyapps.io / Shiny Server / ShinyProxy / \texttt{shiny.collections}. Show that none support real-time multi-user collaboration. ShinyProxy is closest but still uses container-per-user isolation, not shared state.
  
  \item \textbf{Slide: ``Architectural Foundations''} --- Microservices (Fowler 2014) + Kafka event streaming (Kreps 2013) + REST/Plumber (Rodriguez 2021) + Docker for reproducibility (Boettiger 2015). These are the building blocks combined in ShinySwarm.
\end{enumerate}

% ============================================================
\begin{thebibliography}{99}
% ============================================================

\bibitem{baker2016} M. Baker. ``1,500 scientists lift the lid on reproducibility.'' \textit{Nature}, 533, pp.\ 452--454, 2016.

\bibitem{peng2011} R. D. Peng. ``Reproducible Research in Computational Science.'' \textit{Science}, 334(6060), pp.\ 1226--1227, 2011.

\bibitem{rcore2025} R Core Team. \textit{R: A Language and Environment for Statistical Computing.} 2025. \url{https://www.r-project.org/}

\bibitem{chang2022} W. Chang, J. Cheng, J. Allaire, C. Sievert, B. Schloerke, Y. Xie, J. Allen, J. McPherson, A. Dipert, B. Borges. \textit{Shiny: Web Application Framework for R.} 2022. \url{https://shiny.posit.co/}

\bibitem{hartgerink2017} C. Hartgerink. ``Composing reproducible manuscripts using R Markdown.'' 2017.

\bibitem{churchill2001} E. Churchill, D. Snowdon, A. Munro. \textit{Collaborative Virtual Environments: Digital Places and Spaces for Interaction.} Springer, 2001.

\bibitem{aerts2015} A. Aerts. \textit{shinyURL: Save and restore the state of Shiny apps.} 2015. \url{https://github.com/aoles/shinyURL}

\bibitem{shinybookmark} Posit. ``Bookmarking state in Shiny applications.'' \url{https://shiny.posit.co/r/articles/share/bookmarking-state/}

\bibitem{appsilon2017} Appsilon. \textit{shiny.collections: Google Docs-like live collaboration in Shiny.} UseR! 2017. \url{https://github.com/Appsilon/shiny.collections}

\bibitem{rethinkdb} RethinkDB. \url{https://rethinkdb.com/}

\bibitem{richardson_shared_db} C. Richardson. ``Shared Database Pattern.'' \url{https://microservices.io/patterns/data/shared-database.html}

\bibitem{shinyappsio} Posit. shinyapps.io. \url{https://www.shinyapps.io/}

\bibitem{shinyserver} Posit. \textit{Shiny Server Pro Admin Guide.} 2020. \url{https://docs.posit.co/shiny-server/}

\bibitem{shinyproxy} OpenAnalytics. ShinyProxy. 2025. \url{https://shinyproxy.io/}

\bibitem{brokamp2018} C. Brokamp. ``Shiny Microservices Demo.'' useR! 2018 Conference.

\bibitem{rodriguez2021} J. C. Rodriguez. ``Turning our monolithic Shiny app into a microservices-based structure.'' 2021.

\bibitem{plumber2025} B. Schloerke, J. Allen, B. Tremblay, F. van Dunn\'{e}, S. Vandewoude. \textit{Plumber: An API Generator for R.} 2025. \url{https://www.rplumber.io/}

\bibitem{kafka_docs} Apache Software Foundation. \textit{Apache Kafka Documentation.} \url{https://kafka.apache.org/documentation/}

\bibitem{gupta2022} N. Gupta. \textit{rkafka: Using Apache Kafka Messaging Queue Through R.} 2022.

\bibitem{neff2019} P. Neff. ``Connecting to Apache Kafka from R.'' Confluent tutorial, 2019.

\bibitem{rkafka_examples} GitHub examples of rkafka usage. \url{https://github.com/}

\bibitem{neudecker2024} A. Neudecker, INWT Statistics. \textit{kafka: R Client for Apache Kafka.} 2024. \url{https://github.com/INWTlab/r-kafka}

\bibitem{gunjaca_poc} I. Gunjaca. Proof-of-Concept: Kafka + Shiny microservice. \url{https://github.com/IfaPifa/Collaborative-RShiny}

\bibitem{attali2016} D. Attali. ``Persistent data storage in Shiny apps.'' 2016. \url{https://deanattali.com/blog/shiny-persistent-data-storage/}

\bibitem{mongodb} MongoDB Documentation. \url{https://www.mongodb.com/docs/}

\bibitem{docker_docs} Docker Inc. ``What is a Container?'' \url{https://www.docker.com/resources/what-container/}

\end{thebibliography}

\end{document}
